\chapter{Unit 8: Process Management \& Signal Architecture Question Bank}

\section{Multiple Choice Questions (MCQs)}

\mcqitem{8.1}{What is the process identifier (PID) of the systemd / init ancestor process in Linux?}
{\texttt{0}}
{\texttt{1}}
{\texttt{-1}}
{\texttt{100}}

\mcqitem{8.2}{Which system call creates a duplicate child process with an identical address space?}
{\texttt{fork()}}
{\texttt{execve()}}
{\texttt{clone()}}
{\texttt{spawn()}}

\mcqitem{8.3}{What is a Zombie (defunct) process?}
{A process stuck in an infinite CPU loop}
{A terminated process whose exit status has not yet been collected by its parent via wait()}
{A background daemon executing without a controlling terminal}
{A process whose memory has been paged out to swap}

\mcqitem{8.4}{Why cannot a Zombie process be destroyed by sending: kill -9 <PID>?}
{Zombie processes run in Ring 0 with elevated immunity}
{The process has already terminated and released its memory; only its table entry remains}
{The root user lacks permission to signal defunct processes}
{\texttt{kill -9} only targets foreground jobs}

\mcqitem{8.5}{Which POSIX signal cannot be intercepted, blocked, or caught by an application?}
{\texttt{SIGINT (2)}}
{\texttt{SIGTERM (15)}}
{\texttt{SIGKILL (9)}}
{\texttt{SIGHUP (1)}}

\mcqitem{8.6}{Which keyboard shortcut sends the SIGTSTP signal, suspending a foreground process?}
{\texttt{Ctrl+C}}
{\texttt{Ctrl+Z}}
{\texttt{Ctrl+\textbackslash}}
{\texttt{Ctrl+D}}

\mcqitem{8.7}{Which command resumes a stopped process running in the background?}
{\texttt{bg}}
{\texttt{fg}}
{\texttt{resume}}
{\texttt{continue}}

\mcqitem{8.8}{What does the STAT code 'D' in the output of the ps command indicate?}
{Defunct / Zombie state}
{Uninterruptible sleep (typically waiting on disk device I/O)}
{Detached daemon}
{Debug tracing state}

\mcqitem{8.9}{Which command signals processes matching a specified pattern name?}
{\texttt{killall}}
{\texttt{pkill}}
{\texttt{pskill}}
{\texttt{term}}

\mcqitem{8.10}{What is the function of the nohup command?}
{Runs a process with maximum CPU scheduling priority}
{Executes a command immune to SIGHUP terminal disconnect signals}
{Prevents memory allocation errors}
{Disables child subshell creation}

\begin{table}[htbp]
\centering
\caption{\textbf{Unit 8 MCQ Answer Key \& Verification}}
\begin{tabularx}{0.85\linewidth}{c c X}
\toprule
\textbf{Question} & \textbf{Answer} & \textbf{Conceptual Core} \\
\midrule
8.1 & \textbf{(b)} & The \texttt{systemd} or \texttt{init} process is the root of the process tree (PID 1). \\
8.2 & \textbf{(a)} & \texttt{fork()} duplicates the calling process; \texttt{execve()} overwrites it with new code. \\
8.3 & \textbf{(b)} & A zombie has exited, but remains in the process table until parent reaps its status. \\
8.4 & \textbf{(b)} & Zombies are already deceased; signals cannot terminate dead processes. \\
8.5 & \textbf{(c)} & \texttt{SIGKILL (9)} and \texttt{SIGSTOP (19)} cannot be trapped or ignored. \\
8.6 & \textbf{(b)} & \texttt{Ctrl+Z} generates \texttt{SIGTSTP} (signal 20), placing the job into stopped state. \\
8.7 & \textbf{(a)} & \texttt{bg \%N} restarts a stopped job in background mode. \\
8.8 & \textbf{(b)} & State \texttt{D} is uninterruptible sleep, awaiting hardware I/O completion. \\
8.9 & \textbf{(b)} & \texttt{pkill} matches processes by regex pattern and user attributes. \\
8.10 & \textbf{(b)} & \texttt{nohup} catches \texttt{SIGHUP}, ensuring background execution continues after logout. \\
\bottomrule
\end{tabularx}
\end{table}

\section{Technical Viva-Voce Questions \& Conceptual Answers}

\subsection*{Q1. Explain the life cycle of a process from fork() through zombie state and reaping.}
\begin{vivasol}
1. **Creation:** A parent process invokes \texttt{fork()}, creating an identical child process with a unique PID.
2. **Execution:** The child calls \texttt{execve()}, replacing its memory space with a new program image.
3. **Termination:** Upon completing execution, the process calls \texttt{exit(code)}. The kernel frees its virtual memory, open file descriptors, and heap pages.
4. **Zombie Phase:** The kernel retains a minimal entry in the process table holding its PID and exit status code.
5. **Reaping:** The parent invokes \texttt{wait()} or \texttt{waitpid()}, reading the child's exit status. Upon completion of \texttt{wait()}, the kernel removes the process table entry entirely.
\end{vivasol}

\section{Lab Practical Exercises \& Step-by-Step CLI Execution}

\begin{practicalbox}{8.1 --- Terminal Job Control Workflow}
\textbf{Objective:} Launch a long-running simulation job, pause it via keyboard signal, verify background status, and return it to foreground.
\begin{clibox}
$ python3 -c 'import time; [time.sleep(1) for _ in range(300)]'
^Z
[1]+  Stopped                 python3 -c '...'

$ jobs -l
[1]+ 24512 Stopped                 python3 -c '...'

$ bg %1
[1]+ python3 -c '...' &

$ jobs -l
[1]+ 24512 Running                 python3 -c '...' &

$ kill -15 24512
[1]+  Terminated              python3 -c '...'
\end{clibox}
\end{practicalbox}
